\textbf{COLUNI 2026 -- Questão 25}\par\medskip Leia o trecho a seguir.

\textbf{Mortes por envenenamento: Brasil teve mais de 100 casos entre 2020 e 2024}

Entre 2020 e 2024, mais de 100 mortes por envenenamento foram registradas no Brasil. Especialistas estão alarmados com o uso de substâncias que não eram relatadas há tempos, como o arsênio. Esse elemento químico, de número atômico 33, é altamente tóxico e pode estar presente em formas orgânicas e inorgânicas, sendo algumas mais perigosas que outras. Dados do ministério da Saúde mostram que as maiores vítimas desse homicídio são os homens; 70\% são pretos e pardos e 42\% são jovens, entre 21 e 40 anos.\par
Disponível em: https://oglobo.globo.com/brasil/noticia/2025/07/26/mortes-por-envenenamento-brasil-teve-mais-de-100-casos-entre-2020-e-2024.ghtml. Acesso em: 26.jul.2025. (Adaptado).

Com base nas informações da reportagem e nos conhecimentos sobre ciências, assinale a afirmativa CORRETA:\begin{enumerate}
\item O arsênio é um metal alcalino terroso, localizado na família 4 da Tabela Periódica, e reage fortemente com a água.
\item Um átomo neutro de arsênio contém 33 prótons, 33 elétrons e 33 nêutrons.
\item O número atômico do arsênio indica a quantidade de prótons presentes no núcleo de seus átomos.
\item O uso de arsênio como veneno é recente e surgiu a partir de pesquisas modernas em nanotecnologia.\end{enumerate}